% Shared abstract, \input by both main.tex (IEEEtran) and preview.tex (article).
PageRank is a canonical workload for evaluating large-scale graph-processing
systems, yet implementations written for different frameworks are rarely checked
for \emph{numerical equivalence} before their performance is compared. We present
a reproducible, single-codebase study that pairs a dependency-free pure-Python
\emph{reference engine} with five Hadoop-ecosystem implementations of PageRank:
Python \texttt{mrjob} (a local engine and a genuine single-pass MapReduce job),
PySpark (RDD and DataFrame), Hadoop Streaming, Java MapReduce, and Apache Pig.
The reference engine implements the standard, weighted, and personalized variants
with correct dangling-node handling and matches \texttt{networkx.pagerank} to a
maximum per-node error of $2.6\times10^{-12}$; every executable distributed
implementation agrees with it to within $2.4\times10^{-8}$, i.e.\ up to the
eight-decimal output precision, so any timing difference reflects engineering
overhead rather than a different computation. A single configuration-driven
harness measures wall-clock time, peak memory, and accuracy on synthetic
Barab\'asi--Albert graphs. On a single node and at moderate scale the in-memory
engine is fastest, \texttt{mrjob} adds a $1.3\times$--$1.8\times$ overhead, and
the process-per-iteration Hadoop Streaming pipeline is $5\times$--$13\times$
slower while using the least memory. We analyze the architectural causes of these
differences, document a self-normalization property, a single-pass dangling-mass
limitation, and a Spark lineage-growth pitfall, and release the full toolchain so
that cluster-scale evaluation of the remaining frameworks is a single command.
The study argues for a discipline of \emph{validate first, then compare}.
